\documentclass[10pt,a4paper]{amsart}
\usepackage{amsmath,amssymb,amsthm,graphicx}
% Retained preamble for the Stacks Project chapter divisors.tex (Divisors).
% The chapter file pinned-original-source.tex carries only \input{preamble} on its first line, so no
% preamble text lives inside the chapter itself.  The Stacks Project uses one single shared preamble for
% every chapter; the snapshot stages it verbatim in sources/stacks-chow/preamble.tex.  The block below is
% copied mechanically, line for line, from that file's lines 14-18 (the Xy-pic setup) and lines 154-217
% (the native theorem declarations together with the mathematical macros).  Nothing was added, removed or
% redefined.  The wrapper class and the amsmath/amssymb/amsthm/graphicx base packages are supplied by the
% assembler, so the class/\externaldocument part of the shared preamble is not repeated here.
% For commutative diagrams we use Xy-pic
\usepackage[all]{xy}

% We use 2cell for 2-commutative diagrams.
\xyoption{2cell}
% Theorem environments.
%
\theoremstyle{plain}
\newtheorem{theorem}[subsection]{Theorem}
\newtheorem{proposition}[subsection]{Proposition}
\newtheorem{lemma}[subsection]{Lemma}

\theoremstyle{definition}
\newtheorem{definition}[subsection]{Definition}
\newtheorem{example}[subsection]{Example}
\newtheorem{exercise}[subsection]{Exercise}
\newtheorem{situation}[subsection]{Situation}

\theoremstyle{remark}
\newtheorem{remark}[subsection]{Remark}
\newtheorem{remarks}[subsection]{Remarks}

\numberwithin{equation}{subsection}

% Macros
%
\def\lim{\mathop{\mathrm{lim}}\nolimits}
\def\colim{\mathop{\mathrm{colim}}\nolimits}
\def\Spec{\mathop{\mathrm{Spec}}}
\def\Hom{\mathop{\mathrm{Hom}}\nolimits}
\def\Ext{\mathop{\mathrm{Ext}}\nolimits}
\def\SheafHom{\mathop{\mathcal{H}\!\mathit{om}}\nolimits}
\def\SheafExt{\mathop{\mathcal{E}\!\mathit{xt}}\nolimits}
\def\Sch{\mathit{Sch}}
\def\Mor{\mathop{\mathrm{Mor}}\nolimits}
\def\Ob{\mathop{\mathrm{Ob}}\nolimits}
\def\Sh{\mathop{\mathit{Sh}}\nolimits}
\def\NL{\mathop{N\!L}\nolimits}
\def\CH{\mathop{\mathrm{CH}}\nolimits}
\def\proetale{{pro\text{-}\acute{e}tale}}
\def\etale{{\acute{e}tale}}
\def\QCoh{\mathit{QCoh}}
\def\Ker{\mathop{\mathrm{Ker}}}
\def\Im{\mathop{\mathrm{Im}}}
\def\Coker{\mathop{\mathrm{Coker}}}
\def\Coim{\mathop{\mathrm{Coim}}}

% Boxtimes
%
\DeclareMathSymbol{\boxtimes}{\mathbin}{AMSa}{"02}

%
% Macros for moduli stacks/spaces
%
\def\QCohstack{\mathcal{QC}\!\mathit{oh}}
\def\Cohstack{\mathcal{C}\!\mathit{oh}}
\def\Spacesstack{\mathcal{S}\!\mathit{paces}}
\def\Quotfunctor{\mathrm{Quot}}
\def\Hilbfunctor{\mathrm{Hilb}}
\def\Curvesstack{\mathcal{C}\!\mathit{urves}}
\def\Polarizedstack{\mathcal{P}\!\mathit{olarized}}
\def\Complexesstack{\mathcal{C}\!\mathit{omplexes}}
% \Pic is the operator that assigns to X its picard group, usage \Pic(X)
% \Picardstack_{X/B} denotes the Picard stack of X over B
% \Picardfunctor_{X/B} denotes the Picard functor of X over B
\def\Pic{\mathop{\mathrm{Pic}}\nolimits}
\def\Picardstack{\mathcal{P}\!\mathit{ic}}
\def\Picardfunctor{\mathrm{Pic}}
\def\Deformationcategory{\mathcal{D}\!\mathit{ef}}

% --- appended 2026-10-12 (Chief): the project own definitions of the reference/slogan/history
% display environments, copied mechanically from the shared Stacks preamble lines 8-12
% (\usepackage{verbatim} plus the three \newenvironment lines, which the chapter preamble keeps
% verbatim from the project preamble). Nothing added, removed or redefined: these environments
% carry no mathematical content and the project itself typesets them as comments. Without these
% three lines a counted statement whose own text uses them fails with Environment slogan undefined.
\usepackage{verbatim}
\newenvironment{reference}{\comment}{\endcomment}
\newenvironment{slogan}{\comment}{\endcomment}
\newenvironment{history}{\comment}{\endcomment}
\usepackage{cleveref}
\crefname{originallemma}{original lemma}{original lemmas}
\Crefname{originallemma}{Original lemma}{Original lemmas}
\makeatletter
\newlabel{coherent-lemma-coherent-Noetherian}{{01XZ}{source}{}{sourcepointer-1}{}}
\newlabel{coherent-lemma-coherent-Noetherian@cref}{{[originallemma][0][]01XZ}{[1][0][]source}{}{}{}}
\makeatother
\crefname{originallemma}{original lemma}{original lemmas}
\Crefname{originallemma}{Original lemma}{Original lemmas}
\makeatletter
\newlabel{constructions-lemma-base-change-map-proj}{{01N2}{source}{}{sourcepointer-2}{}}
\newlabel{constructions-lemma-base-change-map-proj@cref}{{[originallemma][0][]01N2}{[1][0][]source}{}{}{}}
\makeatother
\crefname{originallemma}{original lemma}{original lemmas}
\Crefname{originallemma}{Original lemma}{Original lemmas}
\makeatletter
\newlabel{lemma-relative-proj-finite-type}{{07ZY}{source}{}{sourcepointer-3}{}}
\newlabel{lemma-relative-proj-finite-type@cref}{{[originallemma][0][]07ZY}{[1][0][]source}{}{}{}}
\makeatother
\crefname{originallemma}{original lemma}{original lemmas}
\Crefname{originallemma}{Original lemma}{Original lemmas}
\makeatletter
\newlabel{morphisms-lemma-noetherian-finite-type-finite-presentation}{{01TX}{source}{}{sourcepointer-4}{}}
\newlabel{morphisms-lemma-noetherian-finite-type-finite-presentation@cref}{{[originallemma][0][]01TX}{[1][0][]source}{}{}{}}
\makeatother
\title{The relative Proj of a finitely presented graded algebra is of finite presentation with finitely presented twists}
\author{The Stacks Project Authors}
\date{stacks-project:divisors.tex:local-snapshot-20261012; GNU FDL 1.2 or later}
\begin{document}
\maketitle
\noindent\textit{Editorial note.} All delivered mathematical text is the author's own, copied from the pinned Stacks Project chapter file pinned-original-source.tex; no statement or proof step was rewritten, reordered or removed.\par\medskip
\noindent\textit{Editorial note.} Following the multi-claim convention for this source, source\_url is the source-level URL https://stacks.math.columbia.edu/ and the source label lemma-relative-proj-finite-presentation is carried in claim\_key and here rather than in a per-claim tag URL.\par\medskip
\noindent\textit{Editorial note.} The preamble is the chapter's shared Stacks Project preamble, retained as sources/stacks-divisors/retained-preamble.tex: a line-for-line excerpt of the shared Stacks preamble staged with this snapshot. No macro was invented or redefined, and this card adds no preamble text of its own.\par\medskip
\noindent\textit{Editorial note.} Statement span 7845-7853 and proof span 7855-7886 come from the halfA pool card operations/pool-cards/stacks-divisors-halfA.json and were re-read in the pinned chapter file before this card was written.\par\medskip
\noindent\textit{Editorial note.} The proof writes the graded presentation with the author's weights on the variables and degrees on the defining polynomials; both are preserved.\par\medskip
\noindent\textit{Editorial note.} 4 cross-references occur inside the counted spans; they are listed in external\_references as navigation pointers only.\par\medskip
\noindent\textit{Allowed standard prerequisites (editorial).}\begin{itemize}
\item The compatibility of relative Proj with base change of the graded algebra, cited from the constructions chapter. Cited at proof lines 7869-7871. Lets the statement be checked over a finitely generated model ring and then transported back.
\item The relative Proj of a finitely generated graded algebra is of finite type and its twists are quasi-coherent modules of finite type, cited from the same chapter. Cited at proof lines 7873-7876. Gives finite type for the model, which is upgraded to finite presentation over a Noetherian base.
\item A finite type morphism of Noetherian schemes is of finite presentation, and a coherent module over a Noetherian scheme is of finite presentation, cited from the morphisms chapter and the cohomology of schemes chapter. Cited at proof lines 7878-7883. Upgrades finite type to finite presentation for both the morphism and the twists over the Noetherian model.
\end{itemize}\medskip
\ifdefined\hypertarget\hypertarget{sourcepointer-1}{}\fi
\noindent\textit{Source reference.} Original lemma 01XZ: focused-seeds/AGP/sources/stacks-snapshot-20261011/tags.txt, source lines 2538--2538. Same-chapter reference used in the retained author text at source line(s) 166, 7884. It is supplied only as a navigation pointer so that the retained text compiles; the counted item delivers the author's own statement and the author's complete proof, and no part of the referenced item is copied in or substituted for the counted argument.\par\medskip
\ifdefined\hypertarget\hypertarget{sourcepointer-2}{}\fi
\noindent\textit{Source reference.} Original lemma 01N2: focused-seeds/AGP/sources/stacks-snapshot-20261011/tags.txt, source lines 2143--2143. Same-chapter reference used in the retained author text at source line(s) 7873, 8832. It is supplied only as a navigation pointer so that the retained text compiles; the counted item delivers the author's own statement and the author's complete proof, and no part of the referenced item is copied in or substituted for the counted argument.\par\medskip
\ifdefined\hypertarget\hypertarget{sourcepointer-3}{}\fi
\noindent\textit{Source reference.} Original lemma 07ZY: focused-seeds/AGP/sources/stacks-snapshot-20261011/tags.txt, source lines 10163--10163. Same-chapter reference used in the retained author text at source line(s) 7787, 7875. It is supplied only as a navigation pointer so that the retained text compiles; the counted item delivers the author's own statement and the author's complete proof, and no part of the referenced item is copied in or substituted for the counted argument.\par\medskip
\ifdefined\hypertarget\hypertarget{sourcepointer-4}{}\fi
\noindent\textit{Source reference.} Original lemma 01TX: focused-seeds/AGP/sources/stacks-snapshot-20261011/tags.txt, source lines 2390--2390. Same-chapter reference used in the retained author text at source line(s) 7882. It is supplied only as a navigation pointer so that the retained text compiles; the counted item delivers the author's own statement and the author's complete proof, and no part of the referenced item is copied in or substituted for the counted argument.\par\medskip
\begin{lemma}
\label{lemma-relative-proj-finite-presentation}
Let $S$ be a scheme. Let $\mathcal{A}$ be a quasi-coherent graded
$\mathcal{O}_S$-algebra. Let
$p : X = \underline{\text{Proj}}_S(\mathcal{A}) \to S$ be the relative
Proj of $\mathcal{A}$. If $\mathcal{A}$ is a finitely presented
$\mathcal{O}_S$-algebra, then $p$ is of finite presentation
and $\mathcal{O}_X(d)$ is an $\mathcal{O}_X$-module of finite presentation.
\end{lemma}

\begin{proof}
Affine locally this reduces to the following statement:
Let $R$ be a ring and let $A$ be a finitely presented graded $R$-algebra.
Then $\text{Proj}(A) \to \Spec(R)$ is of finite presentation
and $\mathcal{O}_{\text{Proj}(A)}(d)$ is a
$\mathcal{O}_{\text{Proj}(A)}$-module of finite presentation.
The finite presentation condition implies we can choose
a presentation
$$
A = R[X_1, \ldots, X_n]/(F_1, \ldots, F_m)
$$
where $R[X_1, \ldots, X_n]$ is a polynomial ring graded by giving
weights $d_i$ to $X_i$ and $F_1, \ldots, F_m$ are homogeneous polynomials
of degree $e_j$. Let $R_0 \subset R$ be the subring generated by
the coefficients of the polynomials $F_1, \ldots, F_m$.
Then we set $A_0 = R_0[X_1, \ldots, X_n]/(F_1, \ldots, F_m)$.
By construction $A = A_0 \otimes_{R_0} R$.
Thus by
Constructions, Lemma \ref{constructions-lemma-base-change-map-proj}
it suffices to prove the result for $X_0 = \text{Proj}(A_0)$ over $R_0$.
By Lemma \ref{lemma-relative-proj-finite-type}
we know $X_0$ is of finite type over $R_0$ and
$\mathcal{O}_{X_0}(d)$ is a quasi-coherent $\mathcal{O}_{X_0}$-module
of finite type.
Since $R_0$ is Noetherian (as a finitely generated $\mathbf{Z}$-algebra)
we see that $X_0$ is of finite presentation over $R_0$
(Morphisms, Lemma
\ref{morphisms-lemma-noetherian-finite-type-finite-presentation})
and $\mathcal{O}_{X_0}(d)$ is of finite presentation by
Cohomology of Schemes, Lemma \ref{coherent-lemma-coherent-Noetherian}.
This finishes the proof.
\end{proof}

\end{document}
